\subsection{Fine-Tuned Transformer Model (BERT-base-uncased)}
In terms of error analysis, the confusion matrix provided reveals several important patterns in this model’s predictions. For example, it performs fairly well on genres such as \bf{Drama} and \bf{Comedy}. This indicates that the dataset contains a large number of instances in those genre categories. This boost in performance compared to other genres is likely because of the high amount of training examples, which allows the model to better distinguishing and learn features for those genres. 
The model appears to struggle with less frequency and more ambiguous genres, such as \bf{Film-Noir}, \bf{Fantasy}, and \bf{Western}. This suggests that class imbalance in the dataset played a significant role in the model’s performance.

There is as notable confusion between genres which used similar language in their descriptions. For example, the model appears to frequently misclassify \bf{Action}, \bf{Crime}, and \bf{Drama} among each other. It has similar misclassification issues between the \bf{Horror} and \bf{Thriller} genres. The model also appears to have a bias toward genres with more training data, such as \bf{Drama}. 

These errors can likely be attributed to a few key factors, such as class imbalances in the dataset, single-label requirements for classification when movie genres and descriptions can be subjective, overlapping vocabulary within different movie genre descriptions.

To mitigate these issues in a future model, a few things could be implemented. One potential mitigation would be to allow for multi-label classification instead of single-label. Using a larger, more diverse, and more balanced dataset would also likely improve these errors and the model’s performance. 
